\documentclass[runningheads]{llncs}
\usepackage{longtable}

\begin{document}
\frontmatter
\setcounter{page}{5}

\chapter*{Preface}
\markboth{Preface}{Preface}
This volume contains the proceedings of the 28th International Conference on Coordina-
tion Models and Languages (COORDINATION 2026), held on June 8–12, 2026, in Urbino, Italy hosted by the
University of Urbino Carlo Bo., as part of the 21st International Federated
Conference on Distributed Computing Techniques (DisCoTec 2026).


Modern information systems rely increasingly on combining concurrent, distributed,
mobile, adaptive, reconfigurable, and heterogeneous components. New models, architectures, languages, and verification techniques are necessary to cope with the complexity induced by the demands of today’s software development. Coordination languages
have emerged as a successful approach, in that they provide abstractions that cleanly
separate behaviour from communication, thereby increasing modularity, simplifying
reasoning, and ultimately enhancing software development. Building on the success of
the previous editions, COORDINATION 2026 provides a well-established forum for the
growing community of researchers interested in models, languages, architectures, and
implementation techniques for coordination.


COORDINATION 2026 solicited contributions in four different categories: (1) regular papers describing thorough and complete research results and experience reports;
(2) short papers describing research in progress, calls for action, or opinions on past
COORDINATION research, on the current state of the art, or on prospects for the years
to come; (3) survey papers describing important results and success stories related to
the topics of COORDINATION; and (4) tool papers describing technological artefacts
in the scope of the research topics of COORDINATION, and requiring the submission
of the associated artefact via a separate procedure.


There were 28 paper submissions distributed over the different categories: 24 regular
papers, and 4 tool papers. The selection of the papers
was entrusted to the Programme Committee (PC), consisting of 32 members from 15
different countries (with gender balance F/M = 12/20). The selection of the papers was
done electronically, in two phases. In the first phase, which lasted three weeks, each
submission was single-blind reviewed by at least three PC members, in some cases with
the help of external reviewers. During the second phase, which lasted four days, the papers were thoroughly discussed. The decision to accept or reject a paper
was based not only on the reviews and scores but also on these in-depth discussions. In
the end, 13 papers were selected to be presented at the conference: 11 regular papers, and 2 tool papers.


The authors of the submitted tool papers were requested to participate in the EAPLS
artefact badging, which instead was optional for all the other paper categories. There
were 13 artefact submissions. The 9 members of the Artefact Evaluation Committee
(AEC), chaired by Gianluca Aguzzi, awarded the Available and Reusable Badge to 3
artefacts, and the Available and Functional Badge to 4 artefacts.

As Programme Chairs, we actively contributed to the selection of the three keynote
speakers of DisCoTec 2026:\begin{itemize}
	\item Prof. Mehdi Dastani (Utrecht University, The Netherlands);
	\item Paolo Romano (NESC-ID, Portugal);
	\item Nathalie Bertrand (Inria, France).
\end{itemize}
We are most grateful to Prof. Dastani for accepting our invitation as the  COORDINATION-related keynote speaker. This volume includes the abstract of his keynote
talk: “Efficient Communication and Coordination in Multiagent Reinforcement Learning”.
As is traditional in DisCoTec, a joint session with the best papers from each main conference was organised. The best paper of COORDINATION 2026 was “X” by X.
We are grateful to all the persons involved in COORDINATION 2026. In particular,
we thank the authors for their submissions, the attendees of the conference for their
participation, the PC members and external reviewers for their work in reviewing sub-
missions and participating in the discussions, the AEC chair Gianluca Aguzzi and
the AEC members for their effort in the evaluation of the artefacts, and the Steering Committee,
chaired by Michele Loreti, for their guidance and support. We are also grateful to the
Organising Committee, chaired by Claudio Antares Mezzina, for their excellent job. We also thank
the providers of the EasyChair Conference Management System, which was of great
help in the submission and reviewing process and in the preparation of the proceedings.
We would also like to acknowledge the prompt and professional support from Springer,
which published these proceedings in printed and electronic volumes as part of their
LNCS book series.
\medskip

\begin{flushright}
\noindent April 2026\hfill
{[Roberto Casadei]}\\
{[Fatemeh Ghassemi]}
\end{flushright}


\chapter*{Organization}
\markboth{Organization}{Organization}
\setlength{\tabcolsep}{0pt}     % remove spacing between columns in the following long tables
\setlength{\LTpre}{0pt}         % remove extra vertical space before the long tables

%\section*{General Chair}
%\begin{longtable}{p{0.3\textwidth}p{0.7\textwidth}}
%Robert Bai          & University of Cambridge, UK
%\end{longtable}
%
%\section*{Program Committee Chairs}
%\begin{longtable}{p{0.3\textwidth}p{0.7\textwidth}}
%Vincent Garc\'{\i}a & Technische Universit\"{a}t Dresden, Germany\\
%Laurent Ho          & Dublin City University, Ireland\\
%Steffen Rao         & Carleton University, Canada\\
%\end{longtable}
%
%\section*{Steering Committee}
%\begin{longtable}{p{0.3\textwidth}p{0.7\textwidth}}
%Michael Agarwal     & KU Leuven, Belgium\\
%John Banerjee       & McMaster University, Canada\\
%Antonio Chan        & Harvard University, USA\\
%Stefano Guo         & Ritsumeikan University, Japan\\
%Mohamed Fang        & University of Porto, Portugal\\
%Alessandro Davis    & University of the Aegean, Greece\\
%Rafael Hoffmann     & Hanyang University, South Korea\\
%Mike L\'{o}pez      & Universidad de Granada, Spain\\
%Marek Keller        & University of Calabria, Italy\\
%Adrian Kwon         & Massachusetts Institute of Technology, USA\\
%Hans Meyer          & University of Kaiserslautern, Germany\\
%Nicholas Pinto      & Siemens AG, Germany\\
%Marcus Reddy        & ETH Zurich, Switzerland\\
%Sandra Ryu          & University of Bonn, Germany\\
%Patricia Silva      & Universidad Aut\'{o}noma de Madrid, Spain\\
%Edward  Smith       & King's College London, UK\\
%Tom\'{a}\v{s} Weber & University of Turku, Finland\\
%Gerhard Wu          & Masaryk University, Czech Republic\\
%Hiroki Yamamoto     & Bielefeld University, Hungary\\
%Werner Yoshida      & University of Catania, Italy\\
%\end{longtable}

\section*{Program Committee}
\begin{longtable}{p{0.3\textwidth}p{0.7\textwidth}}
	
S. Akshay &	IIT Bombay, India\\
Duncan Paul Attard & University of Malt, Malta \\
Giorgio Audrito &	University of Turin, Italy \\
Simon Bliudze &	INRIA, France \\
Valentina Castiglioni &  Eindhoven University of Technology, Netherlands \\
Mariangiola Dezani-Ciancaglini &	University of Turin, Italy\\
Cinzia Di Giusto &	Université Côte d’Azur, France \\
Francisco Ferreira & University of London, UK \\
Silvia Ghilezan & University of Novi Sad, Mathematical Institute SASA, Serbia\\
Susanne Graf & CNRS,VERIMAG, France\\
Ludovic Henrio & CNRS, France; IRIF \\
Thomas T. Hildebrandt & 	University of Copenhagen, Denmark \\
Ping Hou & University of Oxford, UK \\
Eva Kühn & Vienna University of Technology, Austria \\
Ivan Lanese & University of Bologna, Italy \\
Carlos López Pombo &  Universidad Nacional de Río Negro, Argentina; CONICET \\
Michele Loreti &	University of Camerino, Italy \\
Stefano Mariani & University of Modena and Reggio Emilia, Italy \\
Hernán C. Melgratti & University of Buenos Aires, Argentina \\
Fabrizio Montesi & University of Southern Denmark, Denmark \\
J. Garreth Morris & University of Iowa, USA\\
Rumyana Neykova & Brunel University London, UK \\
José Proença & University of Porto, Portugal \\
Violet Ka I Pun & Western Norway University of Applied Sciences, Norway \\
António Ravara & 	NOVA LINCS, Portugal \\
Marjan Sirjani & 	Mälardalen University, Sweden \\
Carolyn L. Talcott & 	SRI International, USA \\
Maurice H. ter Beek & ISTI-CNR, Italy \\
Francesco Tiezzi & University of Florence, Italy \\
Nils Timm & University of Pretoria, South Africa \\	
Emilio Tuosto & Gran Sasso Science Institute, Italy \\
Frank Valencia & École Polytechnique de Paris, France
\end{longtable}

\section*{Artefact Evaluation Committee}
\begin{longtable}{p{0.3\textwidth}p{0.7\textwidth}}
Davide Domini & University of Bologna, Italy \\
Nicolas Farabegoli & University of Bologna, Italy\\
Matteo Magnini & University of Luxembourg, Luxembourg\\
Mário Pereira & NOVA School of Science and Technology, Portugal\\
Marco Quadrini & University of Camerino, Italy\\
Corentin Reuther & University of Namur, Belgium \\
Lorenzo Rossi & University of Camerino, Italy \\
Gerard Tabone & University of Malta, Malta\\
Laura Voinea & University of Glasgow, UK
\end{longtable}

\section*{Additional Reviewers}
\begin{multicols}{2}
\noindent Nicola Del Giudice \\
Masoud Ebrahimi \\
Gabriele Genovese \\
Ramchandra Phawade \\
Jose Ignacio Requeno\\
Volker Stolz \\
\end{multicols}

\tableofcontents

\mainmatter

%Create dummy chapters for the sample ToC

%-----keynote
\setcounter{page}{1}
\addtocontents{toc}{\protect\section*{Keynote}}
\author{Mehdi Dastani}
\title{Efficient Communication and Coordination in Multiagent Reinforcement Learning}
\institute{Utrecht University, Netherlands}
\maketitle
\begin{abstract}
	Multiagent reinforcement learning (MARL) is a powerful framework for learning effective policies in complex multiagent environments. A key challenge in MARL is to learn coordinated policies efficiently in a decentralised manner. Prior work has explored a range of communication and coordination mechanisms to address these challenges. In this talk, I will present our recent contributions towards improving the efficiency and performance of MARL systems through principled communication and coordination mechanisms. First, I introduce a decentralised communication scheduling approach that leverages supervised learning to construct a message estimation model. This model enables individual agents to selectively decide when sharing local information is beneficial, thereby reducing unnecessary communication. Empirical results show that this approach significantly decreases communication overhead while improving overall learning performance. Second, I present two logic-based methods for synthesising multiagent reward machines and action-masking artefacts. These methods provide formal guarantees that the resulting policies are both coordinated and safe, while also improving sample efficiency. 
\end{abstract}
\clearpage


%first heading
\setcounter{page}{1}
\addtocontents{toc}{\protect\section*{Theoretical models and foundations for coordination}}
\author{One Author and Another Author}
\title{This Is a Sample Title}
\maketitle
\clearpage

\setcounter{page}{11}
\author{Some Author}
\title{This Is Another Contribution Title}
\maketitle
\clearpage

\setcounter{page}{21}
\author{First Author and Second Author}
\title{A Third Paper Title}
\maketitle
\clearpage

\setcounter{page}{31}
\addtocontents{toc}{\protect\section*{Software Engineering and Applications for Coordinated Systems}}
\author{Still Another Author}
\title{A Paper About Something}
\maketitle
\clearpage

\setcounter{page}{41}
\author{Same Author}
\title{Last Paper of the Sample ToC}
\maketitle
\clearpage

\setcounter{page}{51}
\phantom{Author Index}
\addcontentsline{toc}{title}{\protect\textbf{Author Index}}
\clearpage

\setcounter{page}{31}
\addtocontents{toc}{\protect\section*{Coordination Platforms, Middleware, and Tools}}
\author{Still Another Author}
\title{A Paper About Something}
\maketitle
\clearpage
\end{document}
